\documentclass[11pt,a4paper]{article}

\usepackage[T1]{fontenc}
\usepackage{lmodern}
\usepackage[margin=0.65in,top=0.5in,bottom=0.5in]{geometry}
\usepackage[hidelinks]{hyperref}
\usepackage{enumitem}
\usepackage{titlesec}

\pagestyle{empty}
\setlength{\parindent}{0pt}
\setlength{\parskip}{0pt}
\setlist[itemize]{leftmargin=1.1em, topsep=0.1em, itemsep=0.05em, parsep=0pt}

\titleformat{\section}{\normalsize\bfseries\scshape}{}{0em}{}[\vspace{0.1em}\titlerule]
\titlespacing*{\section}{0pt}{0.6em}{0.3em}

% \role{Organisation}{Dates}{Title}{Location}
\newcommand{\role}[4]{%
    \textbf{#1} \hfill #2\\
    \textit{#3} \hfill \textit{#4}\par
}

\begin{document}

\begin{center}
    {\LARGE\textbf{Kacper Krasowski}}\\[0.35em]
    Turin, Italy \textbar{}
    \href{mailto:kacperkrasowski2000@gmail.com}{kacperkrasowski2000@gmail.com} \textbar{}
    +48 530 732 802 \textbar{}
    \href{https://kacperkrasowski.github.io/}{kacperkrasowski.github.io}\\[0.15em]
    EU citizen, authorized to work in Italy and Spain
\end{center}

\vspace{0.3em}
Third-year PhD student in Economics building theoretical models of digital platforms, generative AI, and social networks. Experienced in close reading of research as a replicator for the Econometric Society and as a teaching assistant in contract theory and optimization.

\section*{Education}

\role{Collegio Carlo Alberto \& Universitat Aut\`onoma de Barcelona}{Sep 2024 -- Present}{PhD in Economics}{Turin, Italy / Barcelona, Spain}
\vspace{0.3em}
\role{Universitat Aut\`onoma de Barcelona}{Oct 2022 -- Jul 2024}{MRes in Economics, Erasmus Mundus Joint Master Degree}{Barcelona, Spain}
\vspace{0.3em}
\role{Warsaw School of Economics}{Oct 2019 -- Jul 2022}{B.A. in Economics, Quantitative Methods in Economics and Information Systems}{Warsaw, Poland}

\section*{Research}

\textbf{GenAI and the Internet} (2026)\\
Equilibrium model of users choosing between AI and the web and of content creators' entry decisions; shows that GenAI's effect on content supply and consumed quality is non-monotonic in AI capability.

\vspace{0.35em}
\textbf{Impact of Censorship in Social Media} (2024)\\
Friedkin--Johnsen opinion-dynamics model of a threshold censorship policy; characterizes its effects on polarization, disagreement, and internal conflict, and how optimal censorship depends on homophily.

\vspace{0.35em}
\textbf{Social Media and Labor Market Outcomes: The Effect of Facebook on College Graduates' Earnings} (2025), with Luca Marcian\`o\\
Difference-in-differences evidence from Facebook's staggered college rollout on graduates' earnings.

\section*{Research \& Professional Experience}

\role{The Econometric Society}{Mar 2026 -- Present}{Replicator}{}
\begin{itemize}
    \item Independently verify the reproducibility of submitted research by testing authors' data, code, and documentation against the reported results.
\end{itemize}
\vspace{0.3em}
\role{Review of Economic Studies}{Sep 2025 -- Present}{Social Media Manager}{}
\begin{itemize}
    \item Manage the journal's social media accounts and digital communications; build and maintain the REStud Tour websites.
\end{itemize}
\vspace{0.3em}
\role{Shopee Pte. Ltd.}{Mar 2022 -- Sep 2022}{Data Analyst}{}
\begin{itemize}
    \item Built web scrapers in Python and analyzed the collected data in Excel.
\end{itemize}

\section*{Teaching}

\textbf{Teaching Assistant:} Contracts, Information and Incentives in Organizations (UAB, 2026); Optimization (UAB, 2026); Econometrics I (University of Turin, 2025)\\
\textbf{Instructor:} Statistics; Computational Statistics (Barcelona School of Economics, 2026)

\section*{Presentations \& Awards}

\textbf{Presentation:} BSE PhD Jamboree, Barcelona (2026)\\
\textbf{Awards:} UAB IDEA PhD Track Fellowship, Barcelona School of Economics (2023); Social Innovation Award, Social Hackathon (2024)

\section*{Skills}

\textbf{Programming:} Python, R, Stata, MATLAB, SQL, C, Bash, \LaTeX\\
\textbf{Languages:} Polish (native), English (proficient), Spanish (intermediate), Italian (intermediate)

\end{document}
